% Volume I: The Epistemology of Suspicion
% Part I. The Power of Paranoia
% Chapter 4: Paranoia, Skepticism, and Conspiracy
% Spine: proof-spines/ch004-spine.md

\chapter{Paranoia, Skepticism, and Conspiracy}
\label{ch:ch004}


A conspiracy hypothesis is not disqualified merely because it proposes coordination; many institutions are coordinated by definition. The decisive question is instead whether the hypothesis remains open to exclusion. Suspicion becomes epistemically pathological when every possible observation is absorbed as confirmation:
\[
P(E\mid H)\approx P(E'\mid H)
\]
for mutually incompatible observations \(E\) and \(E'\). At that point the hypothesis no longer risks contact with the world, because nothing that happens can count against it.

The chapter therefore distinguishes \textbf{open suspicion} from \textbf{self-sealing suspicion}. Open suspicion specifies possible falsifiers, competing explanations, and thresholds of action. Self-sealing suspicion treats missing evidence as proof of concealment and counterevidence as proof of infiltration. The issue is not whether hidden coordination is imaginable, but whether inquiry preserves a meaningful relation between what is observed and what one is prepared to conclude.
